% Results
\section{Results}

\subsection{A Comprehensive Transcriptomic Atlas of Porcine Development}

To define a molecular standard for porcine development, we constructed a large-scale reference atlas using RNA-seq profiles from the PigGTEX consortium (Fig.~\ref{fig:methodology}a). The dataset comprises 1,924 samples spanning five major tissues (Brain, Liver, Muscle, Lung, and Blood) and covering the entire postnatal lifespan. These five tissues were selected for machine learning analysis based on sample availability and developmental stage representation. Current staging primarily relies on chronological age, which is a poor proxy for biological maturity due to inter-individual variability and asynchronous organ development.

 We addressed this by defining five calibrated life stages based on physiological milestones: Infant (0--20d), Early childhood (21--59d), Pre-pubertal (60--149d), Post-pubertal (150--365d), and Adult ($>365$d) (Fig.~\ref{fig:methodology}c). These stages correspond to distinct physiological transitions: weaning completion (~20d), early growth phase (~60d), puberty onset (~150d), and reproductive maturity (~365d). This provided a ground truth framework for training, allowing us to decouple "biological time" from simple chronological time. The resulting dataset serves as the most comprehensive transcriptomic resource for porcine development to date, overcoming the limitations of previous studies that relied on limited timepoints or morphological snapshots.

\subsection{Accurate Inference of Physiological Developmental Stage}

We implemented a calibrated machine learning framework to infer the biological developmental stage directly from tissue-specific transcriptomes (Fig.~\ref{fig:methodology}b). Models were trained on 70\% of samples with a stratified hold-out test set (30\%) to ensure unbiased evaluation. Test set performance demonstrated robust generalization across tissues, with a mean balanced accuracy of 82.5\% (95\% CI: 79.2--85.8\%) across all tissues (Fig.~\ref{fig:performance}a; see Supplementary Table~S1 for detailed metrics).

Performance varied according to tissue complexity and sample availability. The liver, a highly homogeneous metabolic organ with balanced class representation, achieved the highest accuracy (92.0\%, 95\% CI: 88.5--95.5\%), whereas the brain, with its extreme cellular heterogeneity and limited sample size (n=43), showed lower predictability (63.8\%, 95\% CI: 49.2--78.4\%). Notably, tissues with severe class imbalance (e.g., Muscle: only 5 Adult samples) required reduced classification schemes (4-class for Muscle and Liver, 2-class for Brain, Blood, and Lung) to ensure robust model training (Supplementary Table~S2). Crucially, the system achieved near-perfect ordinal consistency in well-sampled tissues like liver ($\rho = 0.98$) and muscle ($\rho = 0.94$) (Fig.~\ref{fig:performance}c). The confusion matrices (Fig.~\ref{fig:performance}b) reveal that "errors" largely occurred between adjacent stages (e.g., Early vs. Pre-pubertal), reflecting the continuous nature of biological development rather than random misclassification. Unlike static morphological grading, our probabilistic calibration captures these transitional states, resolving the fluid progression of development with high confidence.

These tissue-specific performance differences prompted us to ask: what biological programs underlie the model's predictions? To answer this, we examined the functional annotation of the genes that drive accurate stage inference.

\subsection{The Functional Landscape of Development}

Beyond prediction, our model uncovers the biological logic driving maturation. We performed functional enrichment analysis on the top informative genes for each stage, mapping the global landscape of developmental regulation (Fig.~\ref{fig:biovalidation}).

The resulting functional network reveals that each tissue follows a distinct regulatory framework. In the brain, development is driven by synaptic refinement and metabolic support, highlighted by enrichment of \textit{sarcosine metabolism} (linked to NMDA receptor coagonist availability) and \textit{beta-tubulin binding} (critical for axonal transport). In contrast, the liver prioritizes metabolic autonomy and homeostasis, evidenced by the upregulation of \textit{autophagosome-membrane adaptors}. Muscle development centers on contractile machinery assembly and neuromuscular signaling (e.g., \textit{acetylcholinesterase activity}), while the lung focuses on immune system establishment (\textit{IL-6 signaling regulation}) and RNA metabolic remodeling (\textit{m6A reader activity}). This systems-level view moves beyond simple marker genes to capture the coordinated, tissue-specific logic of physiological maturation.

Notably, feature selection showed moderate stability across random seeds (mean Jaccard similarity: 0.16--0.43 across tissues; Supplementary Fig.~S2b), with 18--60\% of top features appearing in $\geq$80\% of runs. This variability reflects biological gene redundancy rather than model instability: selected features consistently show strong correlation with developmental stage (mean |r| = 0.44, 73.6\% significant; Supplementary Fig.~S2c), indicating that multiple valid biomarker combinations exist within each tissue. Critically, cross-tissue feature overlap was extremely low (0.3\% Jaccard similarity vs. 55\% expected by chance), which is 180-fold lower than random expectation. This profound tissue-specificity is explained by distinct expression profiles across tissues (Supplementary Fig.~S2a) and confirms that each tissue uses unique molecular programs for developmental progression.

\subsection{Evolutionary Conservation with Human Development}

A central question for the pig model is whether its developmental timing biologically mirrors that of humans. We addressed this by performing a direct, quantitative comparison with human skeletal muscle transcriptomes from infant to adult stages \cite{schaiter2024} (Fig.~\ref{fig:crossspecies}). \textbf{We note that this cross-species validation is limited to muscle tissue} due to availability of matched human developmental data; validation in other tissues represents an important future direction.

We found a striking conservation of developmental trajectories in muscle tissue. Using an adaptive thresholding approach to optimize signal-to-noise ratio, we identified 36 genes meeting significance criteria (p < 0.10) in both species. There was a highly significant positive correlation in fold-changes between species (Pearson $R = 0.62$, $p = 5.86 \times 10^{-5}$, 95\% bootstrap CI: 0.40--0.78; Supplementary Fig.~S1), with 89\% (32/36) of shared developmental markers changing in the same direction (Fig.~\ref{fig:crossspecies}a). Bootstrap analysis (1000 iterations) confirmed robustness of this correlation (mean R = 0.60, 95\% CI: 0.40--0.78; Supplementary Fig.~S1b), and sensitivity analysis across p-value thresholds (0.05--0.20) showed consistent positive correlations (R > 0.55) when 15--50 genes were included (Supplementary Fig.~S1a). We identified three conserved functional modules driving this maturation. First, consistent with postnatal neuromuscular maturation, we observed synchronized upregulation of \textit{SORCS2} (pig FC=1.84, human FC=0.95), a regulator of motor neuron innervation, and \textit{KREMEN1}, a modulator of Wnt/beta-catenin signaling. Second, structural refinement was marked by \textit{FGFRL1} (pig FC=1.68, human FC=2.10), implicated in muscle development, and \textit{IGSF1}. Third, metabolic transition was evident in the consistent downregulation of \textit{MSS51} and \textit{PDLIM3} in both species. The suppression of \textit{MSS51}, a negative regulator of mitochondrial oxidative metabolism, suggests a conserved mechanism to "release the brake" on oxidative capacity during postnatal growth. 

These findings provide quantitative evidence that pig muscle development mirrors the human program at the molecular level. However, we acknowledge limitations: the human dataset comprises only 4 infant (7--28 months) and 7 adult (30--56 years) samples, representing a substantial age gap that may affect comparability. Additionally, the correlation analysis required parameter optimization (p-threshold and gene number) to maximize signal-to-noise ratio, though sensitivity analysis confirms robustness across parameter choices. The shared ~90\% directional concordance across 36 developmental markers in muscle tissue suggests that therapeutic interventions targeting these pathways in pigs may translate to human applications, though validation in additional tissues and with larger human cohorts would strengthen this conclusion.